% A worked example: the orbifold C/Z_5 and the section z -> z^2.
% Compile with xelatex (twice, for the references).
\documentclass[11pt,a4paper]{article}

\usepackage[margin=2.6cm]{geometry}
\usepackage{amsmath,amssymb,amsthm}
\usepackage{unicode-math}
\setmainfont{STIXTwoText}[Extension=.otf,UprightFont=*-Regular,BoldFont=*-Bold,
  ItalicFont=*-Italic,BoldItalicFont=*-BoldItalic]
\setmathfont{STIXTwoMath-Regular.otf}
\setmonofont[Scale=0.84]{Menlo}
\usepackage{url}
\usepackage{enumitem}
\usepackage[hidelinks]{hyperref}

\setlength{\emergencystretch}{3em}
\setlist{itemsep=2pt,topsep=4pt}

\theoremstyle{plain}
\newtheorem{theorem}{Theorem}[section]
\newtheorem{proposition}[theorem]{Proposition}
\newtheorem{lemma}[theorem]{Lemma}
\newtheorem{corollary}[theorem]{Corollary}
\theoremstyle{definition}
\newtheorem{definition}[theorem]{Definition}
\newtheorem{recall}[theorem]{From the book}
\theoremstyle{remark}
\newtheorem{remark}[theorem]{Remark}

\newcommand{\C}{\mathbb{C}}
\newcommand{\R}{\mathbb{R}}
\newcommand{\Q}{\mathbb{Q}}
\newcommand{\Z}{\mathbb{Z}}
\newcommand{\N}{\mathbb{N}}
\newcommand{\ZZ}{\Z/5}
\newcommand{\abs}[1]{\lvert #1\rvert}
\newcommand{\card}{\sharp}
\newcommand{\Lin}{\operatorname{Lin}}
\newcommand{\DET}{\operatorname{DET}}
\newcommand{\Or}{\operatorname{Or}}
\newcommand{\supp}{\operatorname{supp}}
\newcommand{\domsupp}{\operatorname{dom-supp}}
\newcommand{\id}{\mathrm{id}}
\newcommand{\cl}{\operatorname{cl}}
\newcommand{\coker}{\operatorname{coker}}

\title{A worked example: the orbifold $\C/\Z_5$ and the section $z\mapsto z^2$}
\author{}
\date{On H.~Hofer, K.~Wysocki, E.~Zehnder, \emph{Polyfold and Fredholm Theory}, Springer, 2021}

\begin{document}
\maketitle

\noindent
This note carries out Theorem 16.6.6 (p.~739) of \emph{Polyfold and Fredholm Theory} for Fredholm
index $0$ on one small polyfold: the orbit space $\C/\Z_5$ of the action of the fifth roots of unity,
with the bundle of weight $2$ and the section $z\mapsto z^2$. Every hypothesis is checked against the
definition in the book, from the ep-groupoid and its strong bundle, through the orientation of the
section, the choice of an admissible sc$^+$-multisection and the canonical orientation of the perturbed
solution set, to the value of the invariant, which is $2/5$. The example is small, but it is a genuine
orbifold case: no single-valued equivariant perturbation is transversal, so the multisections of the
book are needed. A second version adds an infinite-dimensional factor with strictly nested levels, so
that sc-calculus and sc-Fredholm theory enter in earnest; the same conclusions hold there.

The proofs given here are complete for the example itself. The conclusion is obtained by applying
Theorem 16.6.6 to the hypotheses verified here. Where a step is an instance of a general statement, that
statement is given in \S\ref{sec:lemmas}, with a proof or with an outline naming the results of the book
it rests on. Every statement in this note has also been checked in a formalization in Lean, and a version
of the note with references to the formal proofs is \texttt{WorkedExample.pdf} in the same repository,
\url{https://github.com/skydancerosel/polyfold_faithful}.

\paragraph{Conventions.} Book references use the numbering and pages of the 2021 edition, for instance
``Definition 15.3.5 (p.~693)''. Each check first recalls the definition of the book in a few lines and
then verifies it. We identify $\C$ with $\R^2$ as a real vector space; ``smooth'' means $C^\infty$. The
group $G=\ZZ$ is written additively.

%======================================================================
\section{The example}\label{sec:example}

Let $\zeta=e^{2\pi i/5}$ and let $G=\ZZ$ act on $\C$ by $g\cdot z=\zeta^g z$.

\subsection{The ep-groupoid}

\begin{recall}[Definitions 1.1.1, 7.1.3, 7.2.1, 7.4.5]\label{rec:epgroupoid}
An \emph{sc-structure} on a Banach space $E$ is a decreasing sequence $E=E_0\supset E_1\supset\cdots$
of Banach spaces such that every inclusion $E_{m+1}\to E_m$ is compact and $E_\infty=\bigcap_m E_m$ is
dense in every $E_m$ (Definition 1.1.1, p.~4). An \emph{ep-groupoid} (Definition 7.1.3, pp.~288--289)
is a groupoid $(X,\mathbf X)$ whose object set and morphism set are M-polyfolds such that
\begin{enumerate}[label=(\arabic*)]
\item the source and target maps $s,t\colon\mathbf X\to X$ are local sc-diffeomorphisms (étale);
\item the unit map and the inverse map are sc-smooth;
\item every $x\in X$ has an open neighbourhood $V(x)$ such that $t\colon s^{-1}(\cl_X V(x))\to X$ is
  proper;
\item the multiplication $\mathbf X\,{}_s{\times}_t\,\mathbf X\to\mathbf X$ is sc-smooth.
\end{enumerate}
It is \emph{effective} (Definition 7.2.1, p.~307) if for every $x$, every $g\in G_x\setminus\{1_x\}$ and
every open neighbourhood $\mathcal U(g)$ there is $\psi\in\mathcal U(g)$ with $s(\psi)\neq t(\psi)$; an
object with $G_x=\{1_x\}$ is \emph{regular}. It satisfies the \emph{Zhou condition} (Definition 7.4.5,
p.~331) if for every smooth $x$ and $g\in G_x$ whose linearization $Tg\colon T_xX\to T_xX$ is the
identity, $g$ lies in the subgroup $T_x\subset G_x$ of elements acting trivially near $x$.
\end{recall}

The objects of $X=G\ltimes\C$ are the points of $\C$, with the constant sc-structure $\C_m=\C$; the
morphisms are the pairs $(g,z)\in G\times\C$, with source $z$ and target $\zeta^g z$. Composition is
$(h,\zeta^g z)\circ(g,z)=(g+h,z)$, the unit at $z$ is $(0,z)$ and the inverse of $(g,z)$ is
$(-g,\zeta^g z)$.

\begin{proposition}\label{prop:epgroupoid}
$X=G\ltimes\C$ is an ep-groupoid with tame object space and without boundary. It is effective; the
isotropy group of $0$ is $G_0=\{(g,0)\}\cong\ZZ$, every $z\neq0$ is regular, and $X$ satisfies the Zhou
condition. Its orbit space $Z=\abs X=\C/\Z_5$ is paracompact.
\end{proposition}

\begin{proof}
\emph{The sc-structure.} With $\C_m=\C$ for every $m$, the inclusions are the identity of a
finite-dimensional space, hence compact, and $\C_\infty=\C$ is dense in each level. The object space is
the sc-Banach space $\C$ itself, an M-polyfold with the identity as retraction; the identity is a tame
retraction (Definition 2.5.2, p.~53), and $\C$ has no boundary since the partial quadrant is all of $\C$.
The morphism space $G\times\C$ is the disjoint union of the five sheets $\{g\}\times\C$.

(1) On the sheet $\{g\}\times\C$ the source is $(g,z)\mapsto z$ and the target is
$(g,z)\mapsto\zeta^g z$, a real-linear isomorphism of $\C$; both are sc-diffeomorphisms onto $\C$.

(2) The unit $z\mapsto(0,z)$ and the inverse $(g,z)\mapsto(-g,\zeta^g z)$ are smooth on every sheet.

(3) The map $(g,z)\mapsto(g,\zeta^g z)$ is a homeomorphism of $G\times\C$, and the projection
$G\times\C\to\C$ is proper because $G$ is finite. Hence $t$ is proper on $G\times\C$, and in particular
on $s^{-1}(\cl V)=G\times\cl V$ for every open $V$.

(4) The fibred product consists of pairs $((h,\zeta^g z),(g,z))$; it is the disjoint union of copies of
$\C$ indexed by $(g,h)$, and the multiplication $((h,\zeta^g z),(g,z))\mapsto(g+h,z)$ is the identity in
$z$ on each copy.

\emph{Isotropy.} $(g,z)$ is an isotropy morphism iff $\zeta^g z=z$, that is iff $g=0$ or $z=0$. So
$G_0=G\times\{0\}\cong\ZZ$ and $G_z=\{(0,z)\}$ for $z\neq0$.

\emph{Effectiveness.} Only $x=0$ has non-trivial isotropy. Let $g\neq0$ and let $\mathcal U$ be an open
neighbourhood of $(g,0)$. It contains $(g,y)$ for small $y\neq0$, and $s(g,y)=y\neq\zeta^g y=t(g,y)$.

\emph{Zhou condition.} Since $X$ is effective, $T_x=\{1_x\}$ for every $x$. For $x\neq0$ there is
nothing to check. For $x=0$ and $g\in G$, the linearization of $z\mapsto\zeta^g z$ at $0$ is
multiplication by $\zeta^g$ on $T_0\C=\C$, which is the identity only for $g=0$.

\emph{Orbit space.} $Z$ is the quotient of the second countable metrizable space $\C$ by a finite group,
hence second countable and, by the metrizability criterion Corollary 7.3.6 (p.~326), metrizable; so it
is paracompact.
\end{proof}

\subsection{The strong bundle}

\begin{recall}[Section 2.6, Definition 8.3.1]\label{rec:strongbundle}
For sc-Banach spaces $U$ and $F$ the local strong bundle $U\triangleleft F\to U$ carries the double
filtration $(U\triangleleft F)_{m,k}=U_m\oplus F_k$, $0\le k\le m+1$ (Section 2.6, p.~65). A
\emph{strong bundle over the ep-groupoid} $X$ (Definition 8.3.1, p.~375) is a strong bundle
$P\colon W\to X$ over the object M-polyfold together with a strong bundle map
$\mu\colon\mathbf X\,{}_s{\times}_P\,W\to W$ covering $t$, with $P(\mu(\varphi,w))=t(\varphi)$, such that
(1) $\mu$ is a surjective local sc-diffeomorphism and linear on fibres; (2) $\mu(1_x,w)=w$;
(3) $\mu(\psi\circ\gamma,w)=\mu(\psi,\mu(\gamma,w))$.
\end{recall}

Let $W=\C\times\C\to\C$, $(z,w)\mapsto z$, be the trivial bundle $\C\triangleleft\C$ with the constant
levels, so that $W_{m,k}=\C\times\C$ for all $m,k$, and let
\[
  \mu\bigl((g,z),(z,w)\bigr)=(\zeta^g z,\ \zeta^{2g}w),
\]
the action of weight $2$ on the fibre.

\begin{proposition}\label{prop:bundle}
$(P,\mu)$ is a strong bundle over $X$. It has the local regularity property, and
$((P,\mu),\id,\id)$ is a strong polyfold bundle structure on $p\colon Y=\abs W\to Z$.
\end{proposition}

\begin{proof}
The fibred product is $\{((g,z),(z,w))\}\cong G\times\C\times\C$. On the sheet of $g$, $\mu$ is the
real-linear isomorphism $(z,w)\mapsto(\zeta^g z,\zeta^{2g}w)$ of $\C\times\C$. It preserves the (constant)
double filtration, so it is a strong bundle map; it is a local sc-diffeomorphism, already onto on one
sheet, and linear on the fibres. This is (1). Property (2) is the case $g=0$, and (3) is
$\zeta^{2(g+h)}=\zeta^{2h}\zeta^{2g}$.

The local regularity property LR of $X$ (Definition 7.4.19, p.~337) follows from the Zhou condition by
Proposition 7.4.22 (p.~338). Since $X$ is effective, Proposition 8.4.3 (p.~382) gives the local
regularity of the bundle (Definition 8.4.2, p.~382). A strong polyfold bundle structure
(Definition 16.3.1, pp.~721--722) consists of a strong bundle over an ep-groupoid satisfying the Zhou
condition, with paracompact orbit space and the local regularity of the bundle, together with
homeomorphisms of the orbit spaces onto $Y$ and $Z$; here we take $Y=\abs W$, $Z=\abs X$ and the
identities. All requirements were verified above and in Proposition~\ref{prop:epgroupoid}.
\end{proof}

\subsection{The section}

\begin{recall}[Definitions 3.1.4, 3.1.7, 3.6.3, 3.6.4, 12.1.1]\label{rec:fredholm}
Let $K=R(U\triangleleft F)\to O$ be a strong local bundle with retraction $R(u,h)=(r(u),\rho(u)h)$ and
$r(0)=0$. A section germ $(f,0)$ of $K\to O$ has a \emph{filling} $g\colon U\to F$ (Definition 3.1.4,
p.~91) if (1) $f=g$ on $O$ near $0$; (2) $g(y)=\rho(r(y))g(y)$ for $y$ near $0$ implies $y\in O$; and
(3) the linearization at $0$ of $y\mapsto[1-\rho(r(y))]g(y)$, restricted to $\ker Dr(0)$, is an
isomorphism $\ker Dr(0)\to\ker\rho(0)$. A \emph{basic germ} (Definition 3.1.7, p.~94) is an sc-smooth
germ $b\colon(C,0)\to\R^N\oplus W$, with $C=[0,\infty)^k\oplus\R^{n-k}\oplus W$ and $b(0)=0$, whose
$W$-component has the form $w-B(a,w)$ with $B$ a contraction of arbitrarily small constant on every
level near $0$. A germ $(f,x)$ at a smooth point $x$ is an \emph{sc-Fredholm germ}
(Definition 3.6.3, p.~135) if in some strong bundle chart around $x$ it has a filling $g$ and there is a
local sc$^+$-section $s$ with $s(0)=g(0)$ such that $g-s$ is conjugate near $0$ to a basic germ. A
section is \emph{sc-Fredholm} (Definition 3.6.4, p.~135) if it is \emph{regularizing}
($x\in X_m$ and $f(x)\in Y_{m,m+1}$ imply $x\in X_{m+1}$) and its germ at every smooth point is
sc-Fredholm. An \emph{sc-Fredholm section functor} (Definition 12.1.1, p.~566) is a section functor of a
strong bundle over an ep-groupoid that is an sc-Fredholm section of $P$.
\end{recall}

We verify Definition 3.6.3 in a \emph{tame} strong bundle chart, one whose base retraction is tame
(Definition 2.5.2, p.~53), and call the data of Definition 3.6.3 in such a chart \emph{tame germ data}.

Write $q(z)=z^2$ and let
\[
  f(z)=(z,q(z))=(z,z^2).
\]

\begin{proposition}\label{prop:fredholm}
$f$ is an sc-Fredholm section functor of $(P,\mu)$, with tame germ data at every point. Its solution set
is $\{0\}$.
\end{proposition}

\begin{proof}
\emph{Functor.} For a morphism $(g,z)$,
$\mu((g,z),f(z))=(\zeta^g z,\zeta^{2g}z^2)=(\zeta^g z,(\zeta^g z)^2)=f(\zeta^g z)$, so $f$ commutes
with the morphisms.

\emph{Regularity and sc-smoothness.} All levels are equal, so $x\in X_{m+1}$ for every $x$ and $f$ is
regularizing. A smooth map between spaces whose levels are all equal is sc-smooth.

\emph{Germ data at $x_0\in\C$.} Every point is smooth. Take the chart $\Phi(x,v)=(x-x_0,v)$ covering
$x\mapsto x-x_0$, onto the local strong bundle $\C\triangleleft\C$ with $R=\id$, so $r=\id$ and
$\rho\equiv\id$. Its base retraction is tame. The local section is $u\mapsto(u,q(u+x_0))$. We take as
filling the principal part itself, $g(u)=q(u+x_0)$: in Definition 3.1.4, (1) is an equality, (2) holds
because $O=U$, and in (3) both $\ker Dr(0)$ and $\ker\rho(0)$ are $0$. The constant section
$s\equiv q(x_0)$ is an sc$^+$-section, since all levels of the fibre are equal, and $s(0)=g(0)$. The
difference
\[
  (g-s)(u)=q(u+x_0)-q(x_0)
\]
vanishes at $0$ and is a basic germ with $n=N=2$, $k=0$ and $W=0$: its $W$-component is $0=w-B$ with
$w\in W=0$ and $B=0$, and the estimate of Definition 3.1.7 is empty.

\emph{Solutions.} $f(z)=0$ in the fibre means $z^2=0$, that is $z=0$.
\end{proof}

\begin{recall}[Definition 6.2.1]\label{rec:index}
For an sc-Fredholm section $f$ and a smooth point $x$, the \emph{linearizations} of $f$ at $x$ form the
set $\Lin(f,x)=\{(f-s)'(x)+a\}$, where $s$ is an sc$^+$-section with $s(x)=f(x)$ and $a\colon
T_xX\to Y_x$ runs through the sc$^+$-operators (p.~239). All these operators are sc-Fredholm of the
same index, the \emph{index of $f$ at $x$}. Since $s$ is an sc$^+$-section if and only if $-s$ is, the
index is that of $(f+s)'(x)$ for any sc$^+$-section $s$ with $(f+s)(x)=0$.
\end{recall}

\begin{proposition}\label{prop:index}
$f$ has index $0$ at every point.
\end{proposition}

\begin{proof}
At $x_0$ take the constant section $s(z)=(z,-x_0^2)$. Then $(f+s)(x_0)=0$ and $(f+s)'(x_0)$ is the
real-linear map $h\mapsto 2x_0h$ of $\R^2$, whose index is $\dim\ker-\dim\coker=0$ by rank--nullity.
\end{proof}

The section $f$ is a functor, so it induces a section $\abs f$ of $p$ (Definition 16.3.18, p.~729),
which is an sc-Fredholm section of the strong polyfold bundle with representative $(W,f)$
(Definition 16.6.1, p.~736). Since $X$ has no boundary, $\partial Z=\emptyset$.

\subsection{The orientation}\label{sec:orientation}

\begin{recall}[Chapter 6 and Section 12.5]\label{rec:orientation}
For a Fredholm operator $T$ let $\det(T)=\lambda(\ker T)\otimes\lambda^*(\coker T)$, with
$\lambda(V)=\Lambda^{\max}V$ (Definition 6.4.1, p.~246). Over the convex set $\Lin(f,x)$ these lines form
a line bundle $\DET(f,x)$ (Proposition 6.6.1, p.~263), and an \emph{orientation of $(f,x)$} is one of
the two orientations of $\DET(f,x)$ (Definition 6.6.2, p.~263); we write $\Or^f_x$ for this two-element
set. Section 6.6 transports orientations along sc-smooth paths: for a path $\varphi$ one chooses
sc$^+$-sections $s$ with $s(\varphi(t),t)=f(\varphi(t))$ and finitely many $s_1,\dots,s_k$ with
property~(P) of Lemma 6.6.4 (p.~265); the kernels $L_t$ of the stabilized operators form a vector
bundle over $[0,1]$, and the exact sequences (6.27) give isomorphisms $\Phi_{E_t}$ from
$\det((f-s(\cdot,t))'(\varphi(t)))$ to $\lambda(L_t)\otimes\lambda^*(\R^k)$ (pp.~266--267), independent
of the choices by Proposition 6.6.7 (p.~267). This yields the local propagation of
Theorem 6.6.11 (p.~272) and the functor $\Gamma_f$ from the path category $P_X$ (smooth points;
classes of sc-smooth paths up to sc-smooth homotopy with fixed ends) to the orientations,
Theorem 6.6.13 (p.~273). An \emph{orientation} of $f$ (Definition 6.6.16, p.~274) is a continuous
section of $\Or^f\to X_\infty$ for the topology of Theorem 6.6.17 (p.~275). Over an ep-groupoid
(Definition 12.5.4, p.~594) it must in addition satisfy $\varphi_*o_{s(\varphi)}=o_{t(\varphi)}$ for
every smooth morphism $\varphi$, the push-forward being that of Theorem 12.5.3 (p.~594). For a
polyfold section an orientation is one for a representative (Definition 16.6.5, p.~739).
\end{recall}

\paragraph{Positivity at isomorphisms.}
Let $L\in\Lin(f,x)$ be an isomorphism. Then $\ker L=0$ and $\coker L=0$, so
$\det(L)=\lambda(0)\otimes\lambda^*(0)=\R\otimes\R^*$, which is canonically $\R$; equivalently, as on
pp.~701--702 of the book, $\det(L)$ is identified with $\lambda(\ker L)=\Lambda^0(0)=\R$. An element
$o\in\Or^f_x$ restricts to an orientation of the line $\det(L)$.

\begin{definition}\label{def:positive}
$o\in\Or^f_x$ is \emph{positive at the isomorphism $L\in\Lin(f,x)$} if the orientation of
$\lambda(\ker L)=\Lambda^0(0)=\R$ induced by $o$ is the ray of $1$, the empty wedge.
\end{definition}

Of the two elements of $\Or^f_x$ exactly one is positive at a given isomorphism $L$, since the two
induce opposite orientations of $\det(L)$.

For $y\in\C$ let $s_y$ be the constant section $z\mapsto(z,y^2)$. Then $(f-s_y)(y)=0$ and
$(f-s_y)'(y)$ is $h\mapsto2yh$, an isomorphism for $y\neq0$.

\begin{proposition}[Existence]\label{prop:orexists}
$f$ has a continuous orientation over $\C$ (Definition 6.6.16) which is positive at $(f-s_1)'(1)=2$.
\end{proposition}

\begin{proof}
Two sc-smooth paths in $\C$ with the same end points are sc-smoothly homotopic with fixed ends by the
straight-line homotopy, so the path category $P_{\C}$ has at most one morphism between any two points.
Choose $o_0\in\Or^f_0$ and set $o_y=\Gamma_f([\varphi])(o_0)$ for any path $\varphi$ from $0$ to $y$; this
does not depend on $\varphi$. Thus $f$ is orientable (Definition 6.6.14, p.~274), and the section
$y\mapsto o_y$ is continuous for the topology of Theorem 6.6.17. If $o$ is not positive at $2$, the
opposite orientation $y\mapsto -o_y$ (the other element of each $\Or^f_y$) is again continuous and is
positive at $2$.
\end{proof}

\begin{proposition}[Positivity]\label{prop:orpositive}
Let $o$ be a continuous orientation of $f$ positive at $(f-s_1)'(1)$. Then $o$ is positive at
$(f-s_y)'(y)=2y$ for every $y\neq0$.
\end{proposition}

\begin{proof}
Choose $w\in\C$ with $e^w=y$ and let $\gamma(t)=e^{tw}$, a smooth path in $\C\setminus\{0\}$ from $1$ to
$y$. On $\R\times\C$ consider the section $S(t,x)=(x,-\gamma(t)^2)$ of the pulled-back bundle; it is an
sc$^+$-section, being smooth with values in a fibre whose levels are all equal. For every $t$,
$(f+S_t)(\gamma(t))=\gamma(t)^2-\gamma(t)^2=0$, and $(f+S_t)'(\gamma(t))$ is $h\mapsto2\gamma(t)h$, an
isomorphism since $\gamma(t)\neq0$. By Lemma~\ref{lem:path}, $o$ is positive at
$(f+S_1)'(\gamma(1))=(f-s_y)'(y)$ iff it is positive at $(f+S_0)'(\gamma(0))=(f-s_1)'(1)$.
\end{proof}

\begin{proposition}[Invariance]\label{prop:orinvariant}
Every continuous orientation $o$ of $f$ over $\C$ satisfies $\varphi_*o_{s(\varphi)}=o_{t(\varphi)}$ for
every morphism $\varphi$ of $X$, and so is an orientation of the section functor $f$ in the sense of
Definition 12.5.4.
\end{proposition}

\begin{proof}
\emph{At the isotropy of $0$.} For $g\in G$ the push-forward along $(g,0)$ is a permutation $e_g$ of the
two-element set $\Or^f_0$. The push-forward is compatible with composition and is the identity along
units, so $e_g^5=e_{5g}=e_0=\id$. A permutation of a two-element set satisfies $e^2=\id$, hence
$e_g=e_g^5=\id$.

\emph{Along a sheet.} Let $g\in G$. By Theorem 12.5.3 the map $\Or^f\to X_\infty$ is a two-sheeted local
homeomorphism and the local push-forward $\hat\Phi_*$ along the sheet of $(g,z)$ carries continuous
local sections to continuous local sections. Near $t(g,z)$ the sections $o$ and $\hat\Phi_*o$ are two
continuous sections of this double cover over a connected neighbourhood, so they agree either everywhere
or nowhere there. Hence the set of $z\in\C$ for which invariance holds at $(g,z)$ is open and closed. It
contains $0$, and $\C$ is connected, so it is all of $\C$. Lemma~\ref{lem:inv} is the general form of
this argument.
\end{proof}

By Propositions~\ref{prop:orexists}--\ref{prop:orinvariant} there is an orientation $o$ of $f$ in the
sense of Definition 12.5.4 which is positive at $2y$ for every $y\neq0$; we keep it from now on.
Together with the representative $(W,f)$ it is an orientation of the polyfold section $\abs f$
(Definition 16.6.5).

\begin{remark}\label{rem:det}
Over the space $L(\R^2,\R^2)$ of all real-linear maps, the canonical identification $\det(L)=\R$ at
isomorphisms changes sign between the two components of the invertible maps. Hence $o$ is positive
exactly at the isomorphisms of positive determinant, the complex-linear ones among them. This
description is not used below.
\end{remark}

%======================================================================
\section{Multisections}\label{sec:multi}

\subsection{Why multisections are needed}

\begin{recall}[Definitions 13.2.1, 15.1.7, 15.2.1]\label{rec:multisection}
An \emph{sc$^+$-multisection} (Definition 13.2.1, pp.~607--608) is a functor $\Lambda\colon W\to\Q^+$,
so $\Lambda(\mu(\varphi,w))=\Lambda(w)$, such that every $x_0$ has an open neighbourhood $U(x_0)$, with
the natural action of $G_{x_0}$ (Theorem 7.1.13, p.~293), carrying sc$^+$-sections $s_1,\dots,s_k$ and
weights $\sigma_i\in\Q$, $\sigma_i>0$, with $\sum\sigma_i=1$ and
$\Lambda(w)=\sum_{\{i:\,w=s_i(P(w))\}}\sigma_i$ over $U(x_0)$. The \emph{linearization} of
$\Lambda\circ f$ at $x\in\supp(\Lambda\circ f)$ is the formal sum
$T_{(f,\Lambda)}(x)=\sum_{\{i:\,f(x)=s_i(x)\}}\sigma_i\,(f-s_i)'(x)$, independent of the local
structure (Theorem 15.1.1, p.~684; Definition 15.1.7, p.~687). $(f,\Lambda)$ is \emph{transversal at
$x$} if all operators occurring in $T_{(f,\Lambda)}(x)$ are surjective (Definition 15.2.1, p.~688).
\end{recall}

A single-valued sc$^+$-section functor $s$ of $W$ defines the sc$^+$-multisection $\Lambda_s$ with one
branch of weight $1$: $\Lambda_s(w)=1$ if $w=s(P(w))$ and $0$ otherwise.

\begin{proposition}\label{prop:obstruction}
Let $s$ be an sc$^+$-section functor of $W$. Then $s(0)=f(0)$, so $0\in\supp(\Lambda_s\circ f)$, and
$(f,\Lambda_s)$ is not transversal at $0$. Since $X$ has no boundary, $(f,\Lambda_s)$ is therefore not in
general position (Definition~15.3.5, recalled in \S\ref{sec:genpos}).
\end{proposition}

\begin{proof}
An sc$^+$-section of the trivial bundle with constant levels is the graph $s(z)=(z,\psi(z))$ of a smooth
map $\psi\colon\C\to\C$. That $s$ is a functor means $\mu((g,z),s(z))=s(\zeta^g z)$, that is
$\psi(\zeta z)=\zeta^2\psi(z)$. At $z=0$ this gives $\psi(0)=\zeta^2\psi(0)$; since $\zeta^2\neq1$,
$\psi(0)=0=f(0)$, and $\Lambda_s(f(0))=1>0$. The only operator in $T_{(f,\Lambda_s)}(0)$ is
$(f-s)'(0)=Dq(0)-D\psi(0)=-D\psi(0)=:L$, a real-linear map of $\C$. Differentiating
$\psi(\zeta z)=\zeta^2\psi(z)$ at $0$ gives $L(\zeta h)=\zeta^2L(h)$. Write $L(h)=ah+b\bar h$ with
$a,b\in\C$. Then
\[
  L(\zeta h)=a\zeta h+b\zeta^{-1}\bar h,\qquad \zeta^2L(h)=a\zeta^2h+b\zeta^2\bar h,
\]
and comparing coefficients, $a(\zeta-\zeta^2)=0$ and $b(\zeta^{-1}-\zeta^2)=0$. As
$2\not\equiv\pm1\pmod 5$, both factors are non-zero, so $a=b=0$ and $L=0$, which is not onto $\C$.

For the last sentence: with $\partial X=\emptyset$ the reduced tangent space at $0$ is all of $T_0\C$
(\S\ref{sec:genpos}), so general position at $0$ asks exactly that $L$ be onto.
\end{proof}

\subsection{The five-branch multisection}

For $c\in\C$ let $s_j(z)=(z,\zeta^{2j}c)$, $j\in\ZZ$, and
\[
  \Lambda_c(z,w)=\sum_{j\in\ZZ}\tfrac15\,[\,w=\zeta^{2j}c\,]
  =\tfrac15\,\card\{\,j\in\ZZ : w=\zeta^{2j}c\,\}.
\]
Since $2$ is prime to $5$, $j\mapsto\zeta^{2j}$ is a bijection of $\ZZ$ onto the fifth roots of unity;
for $c\neq0$ the values $\zeta^{2j}c$ are distinct, and $\Lambda_c(z,w)=1/5$ if $w^5=c^5$ and $0$
otherwise.

\begin{proposition}\label{prop:multisection}
For every $c\in\C$, $\Lambda_c$ is an sc$^+$-multisection of $(P,\mu)$, with the local section structure
$(s_j)_{j\in\ZZ}$, weights $1/5$, on all of $\C$. The morphisms permute the branches:
$\mu((g,z),s_j(z))=s_{j+g}(\zeta^g z)$. It induces an sc$^+$-multisection $\lambda_c$ of the polyfold
bundle $p$, with $\lambda_c([z,w])=\Lambda_c(z,w)$ (Definition 16.5.1, p.~733).
\end{proposition}

\begin{proof}
\emph{Functor.} $\Lambda_c(\mu((g,z),(z,w)))=\Lambda_c(\zeta^g z,\zeta^{2g}w)$, and
$\zeta^{2g}w=\zeta^{2j}c$ iff $w=\zeta^{2(j-g)}c$; so $j\mapsto j-g$ is a bijection between the two
index sets and the values agree.

\emph{Branches.} Each $s_j$ is a constant graph section, hence an sc$^+$-section (the fibre has constant
levels); the weights are positive with sum $5\cdot\frac15=1$; and
$\Lambda_c(w)=\sum_{\{j:\,w=s_j(P(w))\}}\frac15$ for every $w\in W$, which is the definition. This is a
local representation on the open set $\C$. By Theorem 7.1.13 every $x_0$ has a neighbourhood $U(x_0)$
with the natural $G_{x_0}$-action, and the restriction of the representation to it is a local section
structure in the sense of Definition 13.2.1.

\emph{Permutation.} $\mu((g,z),s_j(z))=(\zeta^g z,\zeta^{2g}\zeta^{2j}c)=s_{j+g}(\zeta^g z)$.

\emph{Descent.} By the functor property $\Lambda_c$ is constant on orbits of $W$, so
$\lambda_c(\abs w)=\Lambda_c(w)$ is well defined on $Y=\abs W$, and $(\lambda_c,\id,W,\Lambda_c)$ is a
representative as required in Definition 16.5.1.
\end{proof}

\begin{proposition}\label{prop:solutions}
$\Lambda_c(f(y))>0$ iff $y^2=\zeta^{2j}c$ for some $j$. If $c\neq0$, this $j$ is unique and
$\Lambda_c(f(y))=1/5$. For $c=r^2$, $r\neq0$, the solutions are the ten points $\pm\zeta^jr$; they form the
two orbits of $r$ and of $-r$, which are distinct, and each solution is regular.
\end{proposition}

\begin{proof}
$f(y)=(y,y^2)$ lies on the branch $s_j$ iff $y^2=\zeta^{2j}c$; for $c\neq0$ the $\zeta^{2j}c$ are
distinct, so exactly one $j$ occurs and the weight is $1/5$. If $c=r^2$, then $y^2=\zeta^{2j}r^2$ iff
$y=\pm\zeta^jr=\zeta^j(\pm r)$. The orbits of $r$ and $-r$ are distinct because $-r=\zeta^gr$ would give
$\zeta^g=-1$, and $-1$ is not a fifth root of unity. The solutions are non-zero, hence regular by
Proposition~\ref{prop:epgroupoid}.
\end{proof}

\subsection{General position}\label{sec:genpos}

\begin{recall}[Definitions 2.5.8, 15.3.5]
For a smooth point $x$ of an M-polyfold the \emph{reduced tangent space} $T^R_xX$ is defined in
Definition 2.5.8 (p.~57); at a point with $d_X(x)=0$, in particular at every point of a polyfold without
boundary, it is all of $T_xX$. $(f,\Lambda)$ is \emph{in general position} (Definition 15.3.5, p.~693) if
for every $x\in\supp(\Lambda\circ f)$ and every operator $L_i$ occurring in
$T_{(f,\Lambda)}(x)=\sum\sigma_iL_i$ the restriction $L_i|T^R_xX\colon T^R_xX\to W_x$ is surjective.
\end{recall}

\begin{proposition}\label{prop:genpos}
For $c\neq0$, $(f,\Lambda_c)$ is in general position, and so is the polyfold pair $(\abs f,\lambda_c)$.
\end{proposition}

\begin{proof}
Let $y\in\supp(\Lambda_c\circ f)$. By Proposition~\ref{prop:solutions}, $f(y)=s_j(y)$ for exactly one
$j$, and $T_{(f,\Lambda_c)}(y)=\frac15(f-s_j)'(y)$. Since $s_j$ is constant, $(f-s_j)'(y)$ is
$h\mapsto2yh$, and $y\neq0$ because $y^2=\zeta^{2j}c\neq0$. So it is onto $\C=W_y$, and
$T^R_y\C=T_y\C$. For the polyfold pair, general position is read in a representative, as are the other
notions of Section 16.6, and $(W,f,\Lambda_c)$ is one.
\end{proof}

\subsection{Auxiliary norm, compactness control, and control of \texorpdfstring{$\lambda_c$}{lambda}}

\begin{recall}[Definitions 8.3.5, 15.3.1, 15.3.3, 16.3.9, 16.6.3]\label{rec:norm}
An \emph{auxiliary norm} (Definition 8.3.5, p.~379; again Definition 12.2.1, p.~570) is a continuous
$N\colon W_{0,1}\to[0,\infty)$ such that (1) $N$ restricts to a complete norm on every fibre; (2) if
$P(h_k)$ converges to $x$ and $N(h_k)\to0$, then $h_k\to0_x$ in $W_{0,1}$; (3) $N(\mu(\varphi,k))=N(k)$.
On a strong polyfold bundle an auxiliary norm is one induced from a representative
(Definition 16.3.9, p.~726). For a saturated open $U\supset S_f$, $(N,U)$ \emph{controls compactness}
(Definition 15.3.1, p.~691) if $\{x\in U: N(f(x))\le1\}$ has an orbit space with compact closure in
$\abs X$; for polyfolds, $(n,U)$ controls compactness if $\cl_Z\{z\in U: n(f(z))\le1\}$ is compact
(Definition 16.6.3, pp.~736--737). An sc$^+$-multisection $\Lambda$ is \emph{$(N,U)$-admissible}
(Definition 15.3.3, p.~692) if $N(\Lambda)\le1$ (Definition 13.2.3, p.~608) and every $x$ with
$(t\Lambda)(f(x))>0$ for some $t\in[-1,1]$ lies in $U$; the text after Definition 15.3.3 notes that
$\domsupp(\Lambda)\subset U$ together with $N(\Lambda)\le1$ suffices.
\end{recall}

Theorem 16.6.6 speaks of an sc$^+$-multisection $\lambda$ \emph{controlled by $(n,U)$}. We read this as
the sufficient condition just quoted, transported to the polyfold:
\[
  n(\lambda)(z)\le1\ \text{for all } z\in Z\qquad\text{and}\qquad \domsupp(\lambda)\subset U,
\]
with $n(\lambda)(z)=\max\{n(y):p(y)=z,\ \lambda(y)>0\}$ as in (16.8) (p.~733) and the domain support of
Definition 16.5.2 (p.~733).

Let $N(z,w)=\abs w$ on $W_{0,1}=\C\times\C$ and let $n$ be the induced function on $Y$. We take $U=Z$.

\begin{proposition}\label{prop:norm}
$N$ is an auxiliary norm of $(P,\mu)$ and $n$ is an auxiliary norm of $p$. The pair $(N,\C)$ controls the
compactness of $f$, and $(n,Z)$ controls the compactness of $\abs f$. For $\abs c\le1$, $\lambda_c$ is
controlled by $(n,Z)$, and $\Lambda_c$ is $(N,\C)$-admissible.
\end{proposition}

\begin{proof}
$N$ is continuous. (1) On the fibre over $z$ it is the modulus of $\C$, a complete norm. (2) If
$z_k\to z$ and $\abs{w_k}\to0$ then $(z_k,w_k)\to(z,0)$. (3)
$N(\mu((g,z),(z,w)))=\abs{\zeta^{2g}w}=\abs w$. By (3), $N$ is constant on orbits and induces $n$ on
$Y$.

$N(f(z))=\abs{z}^2$, so $\{z: N(f(z))\le1\}$ is the closed unit disc. Its image in $Z$ is compact, hence
closed, so it is its own closure in $Z$.

For $\abs c\le1$, every point of $\supp\Lambda_c$ over $z$ is some $s_j(z)$, with
$N(s_j(z))=\abs{\zeta^{2j}c}=\abs c\le1$; so $N(\Lambda_c)\le1$ and $n(\lambda_c)\le1$. The domain
support lies in $Z=U$ and in $\C$, trivially. Admissibility of $\Lambda_c$ is then the remark after
Definition 15.3.3.
\end{proof}

%======================================================================
\section{Theorem 16.6.6 for the example}\label{sec:theorem}

\begin{recall}[Theorem 16.6.6, p.~739]
Let $p\colon Y\to Z$ be a strong polyfold bundle with $\partial Z=\emptyset$ and $f$ an oriented
sc-Fredholm section of index $k$ with compact solution set. Let $n$ be an auxiliary norm and $U$ an open
neighbourhood of $S_f$ such that $(n,U)$ controls compactness. Then for an sc$^+$-multisection $\lambda$
controlled by $(n,U)$ with $(f,\lambda)$ in general position, the branched integrals $\int_{\hat\theta}\tau$
and $\oint_{\hat\theta}\tau$ of a closed sc-differential $k$-form $\tau$ are well defined and independent
of $\lambda$; here $\hat\theta$ is $\theta=\lambda\circ f$ with the orientation inherited from $f$.
\end{recall}

For $k=0$ the inherited orientation is the canonical orientation of Theorem 15.4.1 and
Definition 15.4.2 (p.~702): if near $x$ the multisection has the local structure $(s_i,\sigma_i)$, then
\[
  \hat T(y)=\sum_{\{i:\,f(y)=s_i(y)\}}\sigma_i\cdot\bigl(\ker(f-s_i)'(y),\,o_{i,y}\bigr),
\]
where $o_{i,y}$ is the orientation of $\lambda(\ker(f-s_i)'(y))$ induced by $o_y$ through the
identification $\lambda(\ker L)=\det(L)$ for surjective $L$. The sheets $M_i=\{f=s_i\}$ have tangent
spaces $\ker(f-s_i)'(y)$, of dimension $0$ here, and $\hat\theta=(\theta,\hat T)$ is an oriented branched
suborbifold of dimension $0$ (Definitions 9.1.2, p.~402, and 9.3.2, p.~423). For a $0$-form $\tau$ the
\emph{count} is the branched integral of Definition 9.5.6 (p.~439) for $n=0$,
\[
  \int_{\hat\theta}\tau=\mu^{\hat\theta}_\tau\bigl(\abs{\supp\theta}\bigr),
\]
with the measure of Theorem 9.5.4 (p.~438); its effective version $\oint_{\hat\theta}\tau$ uses the
measure of Theorem 9.5.5 (pp.~438--439). We call an sc$^+$-multisection $\lambda$ an \emph{admissible}
perturbation when it is controlled by $(n,U)$ and $(\abs f,\lambda)$ is in general position, the two
conditions of Theorem 16.6.6.

\begin{theorem}[The hypotheses]\label{thm:hyp}
For $(p,\abs f,o,n,U=Z)$:
\begin{enumerate}[label=(\alph*)]
\item $\partial Z=\emptyset$;
\item the solution set $S_{\abs f}$ is compact;
\item $n$ is an auxiliary norm, and $U=Z$ is an open neighbourhood of $S_{\abs f}$;
\item $(n,Z)$ controls compactness;
\item the representative $(W,f)$ of the orientation has tame germ data at every smooth point, and index
  $0$ at every smooth point.
\end{enumerate}
\end{theorem}

\begin{proof}
(a) Proposition~\ref{prop:epgroupoid}: the object space $\C$ has no boundary. (b) $S_{\abs f}$ is the
single point $[0]$ by Proposition~\ref{prop:fredholm}. (c) Proposition~\ref{prop:norm}; $Z$ is open in
itself. (d) Proposition~\ref{prop:norm}. (e) Propositions~\ref{prop:fredholm} and~\ref{prop:index}.
\end{proof}

\begin{remark}\label{rem:hyp}
Condition (e), together with the orientation $o$ of \S\ref{sec:orientation}, is the form in which we use
the hypothesis of Theorem 16.6.6 that $\abs f$ is an oriented sc-Fredholm section of index $k$: the germ
data of Definition 3.6.3 in tame charts, and the index at every smooth point.
\end{remark}

\begin{corollary}[The conclusion]\label{thm:conclusion}
Any two admissible perturbations of $(\abs f,o)$ whose perturbed solution sets carry the canonical
orientation have the same count and the same effective count of every closed $0$-form on $Z$.
\end{corollary}

\begin{proof}
By Theorem~\ref{thm:hyp} and \S\ref{sec:orientation}, $(p,\abs f,o,n,Z)$ satisfies the hypotheses of
Theorem 16.6.6 (p.~739) with $k=0$, and the corollary is its conclusion.
\end{proof}

We recall why Theorem 16.6.6 holds; the argument is that of Corollary 15.4.4 (p.~703). Two admissible
perturbations $\lambda_0,\lambda_1$ with canonical orientations are joined by an oriented, compact, tame
branched suborbifold of dimension $1$ over $[0,1]\times Z$, built on a saturated neighbourhood of the
compact set where both are controlled, and oriented by Theorem 15.4.1 and Theorem 15.4.3 (p.~702). Its
oriented boundary is $\hat\theta_1-\hat\theta_0$. Stokes' theorem (Theorem 9.5.17, p.~449, and
Theorem 9.5.18, p.~450, for the effective integral) applied to the pull-back of the closed form $\tau$
then gives equal counts.

\begin{theorem}[Admissible perturbations]\label{thm:admissible}
For $0<\abs c\le1$ the multisection $\lambda_c$ is admissible, and $\theta=\lambda_c\circ\abs f$ carries the
canonical orientation $\hat T$ of $(\abs f,o)$, which makes it a compact, tame, oriented branched
suborbifold of dimension $0$. Any two admissible perturbations with the canonical orientation of
$(\abs f,o)$ have the same count and effective count of every closed $0$-form.
\end{theorem}

\begin{proof}
Admissibility is Propositions~\ref{prop:norm} and~\ref{prop:genpos}. Write $\Theta=\Lambda_c\circ f$. By
Proposition~\ref{prop:solutions}, $\supp\Theta$ consists of the ten points $y$ with $y^{10}=c^5$, and
$\Theta(y)=1/5$ there. Near such a $y$ the solution set of $f=s_j$ ($j$ the unique branch through $y$)
is the single point $\{y\}$, because $(f-s_j)'(y)=2y$ is an isomorphism; this is a submanifold of
dimension $0$, and $\Theta=\frac15\cdot\mathbf 1_{\{y\}}$ near $y$ is a local branching structure as in
Definition 9.1.2. The support is finite, hence compact, and tameness (Definition 9.3.6, p.~425) is a
condition at boundary points, of which there are none. The canonical orientation is
$\hat T(y)=\frac15\cdot(0,o_{j,y})$, where $o_{j,y}$ is the orientation of $\lambda(0)=\R$ induced by
$o_y$ at the isomorphism $(f-s_j)'(y)=2y$; by Proposition~\ref{prop:orpositive} it is the positive one.
It is the orientation of Theorem 15.4.1, and the local formula (9.9) of Definition 9.3.2 holds with the
single sheet $\{y\}$. The last sentence is Corollary~\ref{thm:conclusion}.
\end{proof}

\begin{theorem}[The count]\label{thm:count}
For every admissible perturbation with the canonical orientation of $(\abs f,o)$, the count of the
constant form $1$ is
\[
  \int_{\hat\theta}1=\frac25 .
\]
Such perturbations exist: the $\lambda_c$, $0<\abs c\le1$.
\end{theorem}

\begin{proof}
The constant form $1$ is closed (Lemma~\ref{lem:const}). By Theorem~\ref{thm:admissible} the count is
the same for all admissible perturbations with the canonical orientation, so it suffices to compute it
for $\lambda_1$, that is $c=1$.

By Proposition~\ref{prop:solutions} the support of $\Theta=\Lambda_1\circ f$ is the ten points
$\pm\zeta^j$, the orbits of $1$ and of $-1$, which are distinct; each has trivial isotropy, and
$\Theta=1/5$ at each. We apply Lemma~\ref{lem:finite} to the set of representatives $\{1,-1\}$: it lies in
$\supp\Theta$, its two points have different orbits, and their orbits cover the orbit support. So
\[
  \int_{\hat\theta}1=\frac1{\card G_1}\,\varepsilon\bigl(\hat T(1)\bigr)\cdot1
  +\frac1{\card G_{-1}}\,\varepsilon\bigl(\hat T(-1)\bigr)\cdot1,
\]
where $\varepsilon$ is the signed weight. At $y=\pm1$ the only branch through $f(y)$ is $s_0$, and
$(f-s_0)'(y)=2y$ is an isomorphism at which $o$ is positive (Proposition~\ref{prop:orpositive}). By
Lemma~\ref{lem:sign} the signed weight is $\varepsilon(\hat T(y))=\Lambda_1(f(y))=1/5$. With
$\card G_{\pm1}=1$,
\[
  \int_{\hat\theta}1=\frac15+\frac15=\frac25 .\qedhere
\]
\end{proof}

\begin{remark}\label{rem:effective}
The value agrees with the local degree $2$ of $z\mapsto z^2$ at $0$ divided by the order $5$ of the
isotropy group $G_0$. The effective count has the same value, $\oint_{\hat\theta}1=2/5$: since $X$ is
effective and the solutions have trivial isotropy, the effective isotropy groups in Theorem 9.5.5 are
trivial as well, and the computation above applies unchanged.
\end{remark}

%======================================================================
\section{The infinite-dimensional version}\label{sec:inf}

In the example above every level is $\C$. To let the sc-calculus of Part~I, the shifted levels of strong
bundles and sc-Fredholm theory in infinite dimensions enter, we cross it with an sc-Hilbert space whose
levels are strictly nested.

\subsection{The sc-Hilbert space}

Let $H$ be the space of real sequences $x=(x_k)_{k\ge0}$ with the levels
\[
  H_m=\Bigl\{x:\ \abs{x}_m^2=\sum_{k\ge0}(k+1)^{2m}x_k^2<\infty\Bigr\},\qquad m\ge0 .
\]

\begin{proposition}\label{prop:H}
$(H_m)_{m\ge0}$ is an sc-structure on $H_0$ (Definition 1.1.1). Every inclusion $H_{m+1}\to H_m$ is
injective and not onto; for instance $x_k=1/(k+1)$ lies in $H_0$ but not in $H_1$, so not every point of
$H_0$ is smooth.
\end{proposition}

\begin{proof}
\emph{Levels.} The map $x\mapsto((k+1)^mx_k)_k$ is an isometry of $H_m$ onto $\ell^2$, so each $H_m$ is a
Hilbert space, and $H_{m+1}\subset H_m$ with $\abs x_m\le\abs x_{m+1}$.

\emph{Compactness.} In these coordinates, the inclusion $H_{m+1}\to H_m$ becomes the diagonal operator
$D=\operatorname{diag}(1/(k+1))$ on $\ell^2$. Let $D_K$ keep the first $K$ coordinates of $Dy$ and set the
others to $0$. $D_K$ has finite rank and $\lVert D-D_K\rVert=\sup_{k\ge K}1/(k+1)=1/(K+1)\to0$; a norm
limit of finite-rank operators is compact.

\emph{Density.} A finitely supported sequence lies in every $H_m$, hence in $H_\infty$. For $x\in H_m$ the
truncations $x^{(K)}=(x_0,\dots,x_{K-1},0,\dots)$ satisfy
$\abs{x-x^{(K)}}_m^2=\sum_{k\ge K}(k+1)^{2m}x_k^2\to0$, the tail of a convergent series.

\emph{Not onto.} $x_k=(k+1)^{-m-1}$ has $\abs x_m^2=\sum(k+1)^{-2}<\infty$ and
$\abs x_{m+1}^2=\sum1=\infty$; for $m=0$ this is $x_k=1/(k+1)$. Injectivity is clear.
\end{proof}

$H_0$ is a separable Hilbert space, so it has smooth bump functions and is second countable; these give
the bump functions and the paracompactness used below.

\subsection{The polyfold bundle and the section}

The ep-groupoid is $X^\infty=(G\ltimes\C)\times H$, the product of $X$ with $H$ regarded as an
ep-groupoid with only identity morphisms; so the morphisms are $((g,z),e)\colon(z,e)\to(\zeta^gz,e)$. The
strong bundle is $W^\infty=(\C\triangleleft\C)\times(H\triangleleft H)$, with fibre $\C\times H_0$ and
double filtration $\C\times H_k$ over level $m$, $0\le k\le m+1$; the morphisms act by $\zeta^{2g}$ on
the $\C$-factor of the fibre and trivially on the $H$-factor. The section is
\[
  F(z,e)=\bigl((z,z^2),(e,e)\bigr).
\]

\begin{proposition}\label{prop:infbundle}
$X^\infty$ is an ep-groupoid with tame object space $\C\times H$ and no boundary. It is effective, the
isotropy of $(0,e)$ is $\ZZ$ and all other objects are regular, it satisfies the Zhou condition, and its
orbit space $(\C/\Z_5)\times H$ is paracompact. $W^\infty$ is a strong bundle over $X^\infty$ with the
local regularity property, and it gives a strong polyfold bundle $p^\infty$ over
$Z^\infty=(\C/\Z_5)\times H$ with $\partial Z^\infty=\emptyset$. Not every object is smooth.
\end{proposition}

\begin{proof}
The ep-groupoid axioms hold factorwise: on each sheet $\{g\}\times\C\times H$ the source and target are
$(z,e)\mapsto(z,e)$ and $(z,e)\mapsto(\zeta^gz,e)$, and properness, smoothness of the structure maps
and the strong bundle axioms of From the book~\ref{rec:strongbundle} are checked as in
Propositions~\ref{prop:epgroupoid} and~\ref{prop:bundle}, the $H$-factor contributing identities. The
object space $\C\times H_0$ is an sc-Hilbert space, a tame M-polyfold without boundary. The isotropy of
$(z,e)$ is that of $z$ in $X$, so it is $\ZZ$ for $z=0$ and trivial otherwise; effectiveness follows as
in Proposition~\ref{prop:epgroupoid}, using the morphisms $((g,y),e)$, $y\neq0$. For the Zhou condition
at a smooth $(0,e)$: the linearization of $g$ is $(v,h)\mapsto(\zeta^gv,h)$, the identity only for
$g=0$. The orbit space is the product of the metrizable spaces $\C/\Z_5$ and $H_0$. Local regularity
follows from Propositions 7.4.22 and 8.4.3 as before. The point $(0,x)$ with $x_k=1/(k+1)$ is not
smooth, by Proposition~\ref{prop:H}.
\end{proof}

\begin{proposition}\label{prop:inffredholm}
The section $e\mapsto(e,e)$ of $H\triangleleft H$ is sc-Fredholm with tame germ data, and $F$ is an
sc-Fredholm section functor of $W^\infty$ with tame germ data, of index $0$ at every smooth point. Its
solution set is $\{(0,0)\}$.
\end{proposition}

\begin{proof}
\emph{The identity part.} The principal part $e\mapsto e$ is sc-linear, hence sc-smooth. It is
regularizing: if $e\in H_m$ and $(e,e)\in(H\triangleleft H)_{m,m+1}$, the fibre component $e$ lies in
$H_{m+1}$, so $e\in H_{m+1}$. At a smooth point $e_0\in H_\infty$ take the chart
$(e,h)\mapsto(e-e_0,h)$ onto $H\triangleleft H$, whose retraction is the identity and hence tame. The
local principal part is $v\mapsto v+e_0$. As in Proposition~\ref{prop:fredholm} it is its own filling,
since with $R=\id$ the three conditions of Definition 3.1.4 are trivial. The constant section
$s\equiv e_0$ is an sc$^+$-section because $e_0\in H_\infty$: on level $m$ it takes values in $H_{m+1}$,
indeed in every level, and is constant. Then $(g-s)(v)=v$, which is a basic germ with $n=N=0$, $W=H$
and $B=0$; this is the splitting in the proof of Lemma 4.1.11 (p.~179), with finite-dimensional
part $0$.

\emph{The product.} The $\C$-part is Proposition~\ref{prop:fredholm}. The product of two sc-Fredholm
section functors is an sc-Fredholm section functor of the product bundle, as used in the text after
Definition 15.3.3 (p.~692); the tame germ data are the products of the germ data of the factors.

\emph{Index.} At a smooth point $(z_0,e_0)$ let
$s(z,e)=\bigl((z,\,z-z_0-z^2),\,(e,\,-e_0)\bigr)$. Its $\C$-part is a smooth section with values in a
finite-dimensional fibre, and its $H$-part is constant at a smooth point, so $s$ is an sc$^+$-section.
$(F+s)(z,e)=((z,z-z_0),(e,e-e_0))$ vanishes at $(z_0,e_0)$ with linearization the identity of
$\C\times H$, whose index is $0$.

\emph{Solutions.} $F(z,e)=0$ iff $z^2=0$ and $e=0$.
\end{proof}

\begin{proposition}\label{prop:infnorm}
$N^\infty((z,w),(e,h))=\max(\abs w,\abs h_1)$, defined on $W^\infty_{0,1}$ where $h\in H_1$, is an
auxiliary norm, and $(N^\infty,\C\times H)$ controls the compactness of $F$; the induced $n^\infty$
controls the compactness of $\abs F$ with $U=Z^\infty$.
\end{proposition}

\begin{proof}
For $N_H(e,h)=\abs h_1$: (1) on the fibre of $W_{0,1}$, which is $H_1$, it is the complete norm of $H_1$;
(2) if the base points converge and $\abs{h_k}_1\to0$, then $h_k\to0$ in $H_1$; (3) the morphisms act
trivially on $H$. The maximum of auxiliary norms of the two factors is an auxiliary norm of the product,
as in the text after Definition 15.3.3. The set where $N^\infty(F(z,e))\le1$ is
$\{\abs z\le1\}\times\{e\in H_1:\abs e_1\le1\}$. The second factor is the image of the closed unit ball of
$H_1$ under the compact inclusion $H_1\to H_0$, so its closure in $H_0$ is compact. The product of the
two closures is compact, and so is its image in the orbit space.
\end{proof}

\subsection{Orientation, perturbation and count}

\begin{proposition}\label{prop:inforientation}
$F$ has an orientation $o^\infty$ in the sense of Definition 12.5.4 which is positive at the
isomorphisms $(F-F(y,0))'(y,0)=(2y,\id)$ for every $y\neq0$.
\end{proposition}

\begin{proof}
The proofs of Propositions~\ref{prop:orexists}--\ref{prop:orinvariant} apply with the following
changes. \emph{Existence:} two sc-smooth paths in $\C\times H$ with the same ends are sc-smoothly
homotopic with fixed ends by the straight-line homotopy, which is sc-smooth because the space is linear.
\emph{Positivity:} use the path $t\mapsto(e^{tw},0)$ and the sc$^+$-section
$(t,(x,e))\mapsto((x,-e^{2tw}),(e,0))$ of the pulled-back bundle; the linearizations along the path are
$(2e^{tw},\id)$, isomorphisms of $\C\times H$, and Lemma~\ref{lem:path} applies. Normalize at $(1,0)$
as before. \emph{Invariance:} for a morphism $((g,z),e)$ with smooth source and target, the family
$t\mapsto((g,tz),e)$, $t\in[0,1]$, is a connected family of morphisms with smooth sources and targets.
By Lemma~\ref{lem:inv} invariance at $t=1$ is equivalent to invariance at $t=0$, which holds by the
two-element argument of Proposition~\ref{prop:orinvariant}. Definition 12.5.4 asks for invariance only
at smooth morphisms, so morphisms with non-smooth ends impose nothing.
\end{proof}

For $c\in\C$ let $\Lambda^\infty_c$ be the multisection with the five branches
$s^\infty_j(z,e)=((z,\zeta^{2j}c),(e,0))$, each of weight $1/5$; equivalently,
$\Lambda^\infty_c((z,w),(e,h))=\Lambda_c(z,w)$ if $h=0$ and $0$ otherwise.

\begin{theorem}\label{thm:inf}
For the bundle $p^\infty$ and the section $\abs F$ with the orientation $o^\infty$ and $(n^\infty,Z^\infty)$:
\begin{enumerate}[label=(\roman*)]
\item every hypothesis of Theorem 16.6.6 for index $0$ holds, in the form of Theorem~\ref{thm:hyp}, and
  so does its conclusion, in the form of Corollary~\ref{thm:conclusion};
\item for $0<\abs c\le1$ the multisection $\lambda^\infty_c$ induced by $\Lambda^\infty_c$ is admissible and
  its perturbed solution set carries the canonical orientation;
\item every admissible perturbation with the canonical orientation of $(\abs F,o^\infty)$ has count
  $\int_{\hat\theta}1=2/5$.
\end{enumerate}
\end{theorem}

\begin{proof}
(i) $\partial Z^\infty=\emptyset$ and the auxiliary norm are Propositions~\ref{prop:infbundle}
and~\ref{prop:infnorm}; the solution set $\{[(0,0)]\}$ is compact (Proposition~\ref{prop:inffredholm});
compactness control is Proposition~\ref{prop:infnorm}; tame germ data and index $0$ are
Proposition~\ref{prop:inffredholm}; the orientation is Proposition~\ref{prop:inforientation}. The
conclusion is then Theorem 16.6.6, as in Corollary~\ref{thm:conclusion}.

(ii) \emph{Multisection.} As in Proposition~\ref{prop:multisection}: $\Lambda^\infty_c$ is a functor
because the morphisms permute the branches, $((g,z),e)$ carrying $s^\infty_j(z,e)$ to
$s^\infty_{j+g}(\zeta^gz,e)$; the branches are sc$^+$-sections, their $\C$-parts being constant in a
finite-dimensional fibre and their $H$-parts zero; the weights sum to $1$.
\emph{Solutions.} $F(y,e)=s^\infty_j(y,e)$ iff $y^2=\zeta^{2j}c$ and $e=0$. So the solutions are the
$(y,0)$ with $y$ as in Proposition~\ref{prop:solutions}, each on exactly one branch, with weight $1/5$.
\emph{General position.} At $(y,0)$, $(F-s^\infty_j)'(y,0)$ is $(h,k)\mapsto(2yh,k)$, an isomorphism of
$\C\times H$ since $y\neq0$; there is no boundary. \emph{Control.} Every branch has
$N^\infty=\max(\abs c,0)\le1$, and the domain support lies in $Z^\infty$. The canonical orientation is
then obtained as in Theorem~\ref{thm:admissible}.

(iii) By (i) it suffices to compute with $c=1$. The solutions are the ten points $(\pm\zeta^j,0)$,
forming the orbits of $(1,0)$ and $(-1,0)$, which are distinct and regular. At each, the only branch is
$s^\infty_0$ and $o^\infty$ is positive at $(2y,\id)$ (Proposition~\ref{prop:inforientation}), so by
Lemma~\ref{lem:sign} the signed weight is $1/5$. The constant form $1$ is closed (Lemma~\ref{lem:const}),
and Lemma~\ref{lem:finite} with representatives $(1,0),(-1,0)$ gives $\frac15+\frac15=\frac25$.
\end{proof}

%======================================================================
\section{General lemmas}\label{sec:lemmas}

The following statements hold for any strong bundle over an ep-groupoid, or any polyfold, and are used
above.

\begin{lemma}[Counts over finitely many orbits]\label{lem:finite}
Let $\hat\Theta=(\Theta,\hat T)$ be an oriented, tame, compact branched ep$^+$-subgroupoid of
dimension $0$, and let $s$ be a finite subset of $\supp\Theta$ whose points have pairwise different
orbits and whose orbits cover the orbit support $\abs{\supp\Theta}$. Then for every $0$-form $\omega$,
\[
  \int_{\hat\Theta}\omega=\sum_{x\in s}\frac1{\card G_x}\,\varepsilon\bigl(\hat T(x)\bigr)\,\omega(x),
\]
where $\varepsilon(\hat T(x))$, the \emph{signed weight}, is the sum of the weights of the oriented
tangent spaces in $\hat T(x)$, each with the sign of its orientation.
\end{lemma}

\begin{proof}
By Definition 9.5.6, $\int_{\hat\Theta}\omega=\mu(\abs{\supp\Theta})$ for the signed measure
$\mu=\mu^{\hat\Theta}_\omega$ of Theorem 9.5.4. The orbit support is the finite disjoint union of the
singletons $\{\abs x\}$, $x\in s$, so by additivity $\int_{\hat\Theta}\omega=\sum_{x\in s}\mu(\{\abs x\})$.
Let $x\in s$ and choose an oriented local branching structure $(M_i,o_i,\sigma_i)_{i\in I}$ on a neighbourhood
$U(x)$ with the natural $G_x$-action. Each $M_i$ is a submanifold of dimension $0$; after shrinking
$U(x)$ (keeping it $G_x$-invariant) we may assume that every $M_i$ meeting $\{x\}$ is $\{x\}$ and that
the other sheets are empty. Theorem 9.5.4 with $K=\{\abs x\}$ gives
\[
  \mu(\{\abs x\})=\frac1{\card G_x}\sum_{i\in I}\sigma_i\int_{K_i}\omega|(M_i,o_i).
\]
The integral of a $0$-form over the oriented point $(\{x\},o_i)$ is $\varepsilon_i\,\omega(x)$, where
$\varepsilon_i=\pm1$ is the sign of $o_i$. By (9.9) in Definition 9.3.2,
$\hat T(x)=\sum_i\sigma_i\,T_x(M_i,o_i)$, so $\sum_i\sigma_i\varepsilon_i=\varepsilon(\hat T(x))$.
\end{proof}

\begin{lemma}[$d$ of constants]\label{lem:const}
On every M-polyfold the constant $0$-forms are closed: $d\,c=0$. The same holds for the sc-differential
forms on an ep-groupoid.
\end{lemma}

\begin{proof}
A constant $0$-form is the pull-back, by the constant map to a point $\R^0$, of a $0$-form on $\R^0$.
The exterior derivative of sc-differential forms (Section 4.4, p.~195) commutes with pull-back by
sc-smooth maps, and every $1$-form on $\R^0$ vanishes. On an ep-groupoid the forms are the forms on the
object M-polyfold that are invariant under the morphisms; a constant form is invariant, and $d$ is
computed on the object space.
\end{proof}

\begin{lemma}[The sign in dimension $0$]\label{lem:sign}
In the setting of Theorem 15.4.1, let $y$ lie on a sheet $M_i=\{f=s_i\}$ of dimension $0$. The oriented
tangent space $(T_yM_i,o_{i,y})$ is positive if and only if $o_y$ is positive at the isomorphism
$(f-s_i)'(y)$ (Definition~\ref{def:positive}). If $o_y$ is positive at every $(f-s_i)'(y)$ with
$f(y)=s_i(y)$, then $\varepsilon(\hat T(y))=\Lambda(f(y))$.
\end{lemma}

\begin{proof}
Since $T_yM_i=\ker(f-s_i)'(y)$ has dimension $0$ and $(f-s_i)'(y)$ is onto, it is an isomorphism. By
Theorem 15.4.1 the orientation $o_{i,y}$ of $T_yM_i$ is the orientation of
$\lambda(\ker(f-s_i)'(y))=\Lambda^0(0)=\R$ that corresponds to $o_y$ under
$\lambda(\ker L)=\lambda(\ker L)\otimes\R^*=\det(L)$. The oriented point is positive iff this orientation
is the ray of the empty wedge $1$, which is Definition~\ref{def:positive}. For the second statement,
Theorem 15.4.1 gives $\hat T(y)=\sum_{\{i:\,f(y)=s_i(y)\}}\sigma_i\,(T_yM_i,o_{i,y})$, all signs are $+1$,
and $\sum_{\{i:\,f(y)=s_i(y)\}}\sigma_i=\Lambda(f(y))$ by the local representation of
Definition 13.2.1.
\end{proof}

\begin{lemma}[Positivity along paths]\label{lem:path}
Let $X$ be a tame M-polyfold admitting sc-smooth bump functions, $f$ an sc-Fredholm section with tame
germ data of a strong bundle in which every vector over a smooth point lies on an sc$^+$-section, and
$o$ a continuous orientation of $f$ over an open set $U'$. Let $\gamma$ be an sc-smooth path with
$\gamma([0,1])\subset U'$, and $S$ an sc$^+$-section of the pull-back of the bundle to $\R\times X$ such
that $(f+S_q)(\gamma(q))=0$ for $q$ near $[0,1]$ and $(f+S_q)'(\gamma(q))$ is an isomorphism for
$q\in[0,1]$. Then $o$ is positive at $(f+S_0)'(\gamma(0))$ if and only if it is positive at
$(f+S_1)'(\gamma(1))$.
\end{lemma}

\begin{proof}[Outline]
It suffices to show that positivity at $(f+S_q)'(\gamma(q))$ is locally constant in $q\in[0,1]$; then it
is constant, $[0,1]$ being connected. By Theorem 6.6.11 every point $\gamma(q_0)$ has a neighbourhood in
which orientations are propagated along paths by the construction recalled in
From the book~\ref{rec:orientation}, and continuity of $o$ means that $o$ is carried along by this propagation.
Restrict to a subinterval $J\ni q_0$ with $\gamma(J)$ in such a neighbourhood. Use the sections
$S_q$ in the role of $s(\cdot,t)$, and take $k=0$ stabilizing sections: property~(P) of Lemma 6.6.4 asks
for surjectivity, which holds because the operators $(f+S_q)'(\gamma(q))$ are isomorphisms; their kernels are
$0$, so the condition at the boundary is empty. The kernel bundle is then $L_q=0$, the exact sequence (6.27)
is $0\to0\to0\to0\to0$, and $\Phi_{E_q}$ is the canonical identification of $\det((f+S_q)'(\gamma(q)))$
with $\lambda(0)=\R$. The propagation is the continuation of orientations in the line bundle
$\lambda(L)=[0,1]\times\R$, which keeps the ray of the empty wedge $1$. So the orientation transported
from $q_0$ to $q\in J$ is positive iff it was positive at $q_0$, and by continuity of $o$ it is
$o_{\gamma(q)}$.
\end{proof}

\begin{lemma}[Invariance of continuous orientations]\label{lem:inv}
Let $o$ be a continuous orientation of an sc-Fredholm section functor with tame germ data over all of
$X$, whose objects need not all be smooth. Along a continuous family $\gamma\colon T\to\mathbf X$ of
morphisms with $T$ connected and with smooth sources and targets, the condition
$\gamma(t)_*o_{s(\gamma(t))}=o_{t(\gamma(t))}$ either holds for every $t$ or for none.
\end{lemma}

\begin{proof}
Let $\sigma$ be a morphism with smooth source and target and let $\hat\Phi$ be the local push-forward
along the sheet of $\sigma$ (Theorem 12.5.3; p.~594 for the formula). By Theorem 12.5.3, $\hat\Phi_*o$ and
$o$ are continuous sections of the two-sheeted local homeomorphism $\Or^f\to X_\infty$ near
$t(\sigma)$. On a connected neighbourhood of $t(\sigma)$ in $X_\infty$, which exists since $X_\infty$ is
locally connected, two continuous sections of a double cover agree everywhere or nowhere. So for every
$\sigma'$ near $\sigma$ on the same sheet, with smooth source and target, invariance at $\sigma'$ is
equivalent to invariance at $\sigma$. The set of $t\in T$ at which invariance holds at $\gamma(t)$ is
therefore open, and so is its complement; $T$ is connected.
\end{proof}

%======================================================================
\section{Remarks on scope}

\begin{enumerate}[label=(\arabic*)]
\item The first example is $\C$ with all levels equal, so its sc-structure is constant. The version of
  \S\ref{sec:inf} adds an infinite-dimensional factor with strictly nested levels, on which the section is
  the identity.
\item The value $2/5$ is computed for the orientation $o$ normalized to be positive at $(f-s_1)'(1)=2$,
  hence at all $2y$, $y\neq0$: by Remark~\ref{rem:det}, the orientation positive at the complex-linear
  isomorphisms. Theorem~\ref{thm:hyp}, and with it Corollary~\ref{thm:conclusion}, holds for either
  orientation of $\abs f$, since condition (e) concerns only the representative $(W,f)$.
\item The effective count is the same here, $\oint_{\hat\theta}1=2/5$, because $X$ is effective and the
  solutions have trivial isotropy (Remark~\ref{rem:effective}).
\item ``Controlled by $(n,U)$'' in Theorem 16.6.6 is read as $n(\lambda)\le1$ and
  $\domsupp(\lambda)\subset U$, the sufficient condition for admissibility stated after
  Definition 15.3.3.
\end{enumerate}

\end{document}
